% VSR-PathRAG: A Verification--Scoring--Repair Framework for Trustworthy Graph-Based RAG
% Target: IEEE TKDE / Information Processing & Management
%
% Compile with: pdflatex vsr_pathrag.tex
%
\documentclass[journal,onecolumn,11pt]{IEEEtran}

\usepackage{amsmath,amssymb,amsfonts}
\usepackage{amsthm}
\usepackage{algorithmic}
\usepackage{algorithm}
\usepackage{graphicx}
\usepackage{textcomp}
\usepackage{xcolor}
\usepackage{booktabs}
\usepackage{multirow}
\usepackage{hyperref}
\usepackage{url}
\usepackage{cite}

\newtheorem{definition}{Definition}
\newtheorem{theorem}{Theorem}

\def\BibTeX{{\rm B\kern-.05em{\sc i\kern-.025em b}\kern-.08em
 T\kern-.1667em\lower.7ex\hbox{E}\kern-.125emX}}

\begin{document}

\title{VSR-PathRAG: A Verification--Scoring--Repair Framework for Trustworthy Graph-Based Retrieval-Augmented Generation}

\author{
  Author Name \\
  \textit{Affiliation} \\
  \texttt{email@example.com}
}

\maketitle

\begin{abstract}
Retrieval-Augmented Generation (RAG) systems that ground large language model (LLM) outputs in external knowledge suffer from two fundamental trust deficits: (1) \textit{evidence insufficiency}---the retrieved evidence may not logically entail the generated answer, and (2) \textit{claim hallucination}---unsupported or contradicted claims can survive into the final output. Existing GraphRAG approaches retrieve relevant subgraphs but lack a unified runtime contract for deciding whether the evidence is sufficient for the answer's obligations and claims. We propose \textbf{VSR-PathRAG}, an evidence-contract-guided Verification--Scoring--Repair framework. The contract makes question obligations, claim support, evidence necessity, and submission risk explicit; MPS, CPN, and RCS are its selection, intervention, and decision primitives rather than three unrelated contributions. On UltraDomain Mix130, strict ablation over seven configurations shows that the rescue-enhanced VSR layer improves F1 from 0.2206 without repair to 0.2878 while reducing abstention. Under the frozen SAGE v33-preservation instantiation, locked unseen75 plus Repeat1--8 provide 326 zero-overlap holdout examples with complete lexical, same-rubric judge, and pairwise evaluation: SAGE v33 reaches weighted F1 0.3164 versus 0.2972 for both Naive RAG and HyDE, judge overall 3.555 versus 2.960/2.942, and average pairwise target win rates of 59.66\% versus HyDE and 61.50\% versus Naive. A focused mechanism audit further shows that the current contract-gap score identifies harder rows but is not yet an estimate of escalation benefit, making selective upgrade prediction the central next step.
\end{abstract}

\begin{IEEEkeywords}
Retrieval-Augmented Generation, Knowledge Graph, Evidence Verification, Counterfactual Reasoning, Risk Calibration, PAC-Bayesian Bounds
\end{IEEEkeywords}


% ══════════════════════════════════════════════════════════════
\section{Introduction}
% ══════════════════════════════════════════════════════════════

\subsection{Motivation}

Large Language Models (LLMs) generate fluent but potentially ungrounded text. Retrieval-Augmented Generation (RAG) mitigates hallucination by conditioning generation on externally retrieved evidence. However, three critical gaps remain:

\textbf{(1) No proof structure.} Standard RAG retrieves top-$k$ passages and concatenates them into the prompt. There is no formal notion of whether the retrieved set \textit{proves} the answer's claims. This is analogous to presenting evidence without making a legal case.

\textbf{(2) No necessity analysis.} Even when evidence supports the answer, some evidence paths may be \textit{replaceable}---removing them does not affect the answer's correctness. Without identifying \textit{necessary} evidence, the system cannot detect fragile proofs or prioritize retrieval budget.

\textbf{(3) No risk gating.} RAG systems submit answers regardless of evidence quality. When the evidence is insufficient or contradictory, the system should \textit{abstain} or \textit{repair} rather than fabricate.

\subsection{Contributions}

We propose VSR-PathRAG, which addresses these gaps through one evidence-contract problem:

\begin{enumerate}
\item \textbf{Evidence Contract}: We represent a question as obligations and an answer as atomic claims, then compute a contract state describing coverage, support, contradiction, and answer-shape deficits.

\item \textbf{Deficit-Guided Escalation}: MPS and CPN provide the evidence-selection and necessity primitives used to decide whether verification or rescue is worth the additional cost.

\item \textbf{Controlled Repair and Submission}: RCS provides the final decision layer, while deterministic preservation guards prevent a repair from being accepted when it loses supported content without compensating evidence.

\item \textbf{Mechanism-Focused Evaluation}: We report a complete UltraDomain Mix130 ablation, external retrieval-baseline comparisons, representative graph-RAG supplements, locked zero-overlap holdouts with 326/326 complete rows, and an explicit audit of whether contract deficit predicts escalation benefit.
\end{enumerate}


% ══════════════════════════════════════════════════════════════
\section{Problem Formulation}
% ══════════════════════════════════════════════════════════════

\subsection{Notation}

Let $q$ denote a user question, $\mathcal{G} = (V, E)$ a knowledge graph, $\mathcal{E} = \{e_1, \ldots, e_N\}$ the retrieved evidence set, $\mathcal{O} = \{o_1, \ldots, o_P\}$ the question obligations (required answer facets), $\mathcal{C} = \{c_1, \ldots, c_M\}$ the atomic claims extracted from the generated answer, and $\Phi: \mathcal{E} \times \mathcal{C} \to \{+1, -1, 0\}$ the signed claim-evidence graph.

\subsection{Formal Definitions}

\begin{definition}[Answer Trustworthiness]
An answer $a$ with claims $\mathcal{C}$ is \textit{trustworthy} with respect to evidence $\mathcal{E}$ if:
\begin{equation}
\forall c_i \in \mathcal{C}: \exists e_j \in \mathcal{E} \text{ s.t. } \Phi(e_j, c_i) = +1 \text{ and } \nexists e_k \in \mathcal{E} \text{ s.t. } \Phi(e_k, c_i) = -1.
\end{equation}
\end{definition}

\begin{definition}[Proof Subgraph]
A subset $S \subseteq \mathcal{E}$ is a \textit{proof subgraph} for answer $a$ if:
\begin{equation}
\operatorname{cov}_{\mathcal{O}}(S) \geq \tau_O \quad \text{and} \quad \operatorname{cov}_{\mathcal{C}}(S) \geq \tau_C,
\end{equation}
where $\tau_O, \tau_C$ are coverage thresholds. $S$ is \textit{minimal} if no proper subset $S' \subsetneq S$ satisfies both constraints.
\end{definition}

\begin{definition}[Counterfactual Necessity]
Evidence $e_j \in S$ is \textit{necessary} if removing it causes either $\operatorname{cov}_{\mathcal{C}}(S \setminus \{e_j\}) < \tau_C$ or $\operatorname{cov}_{\mathcal{O}}(S \setminus \{e_j\}) < \tau_O$.
\end{definition}

\begin{definition}[Calibrated Risk]
\label{def:risk}
The calibrated risk of submitting answer $a$ is:
\begin{equation}
R(a, S) = \sum_{i=1}^{5} w_i \cdot r_i, \quad \sum_{i=1}^{5} w_i = 1,
\end{equation}
where each $r_i \in [0,1]$ is a normalized risk signal: audit prior, unsupported claim ratio, proof gap, necessity gap, and contradiction risk.
\end{definition}

\begin{definition}[Evidence Contract Deficit]
For question $q$, answer $a$, and evidence set $\mathcal{E}$, the evidence-contract deficit is
\begin{equation}
D(q,a,\mathcal{E}) = w_O(1-\operatorname{cov}_{\mathcal{O}}) + w_C(1-\operatorname{cov}_{\mathcal{C}}) + w_T(1-\operatorname{cov}_{\text{terms}}) + w_S g_{\text{shape}},
\end{equation}
where the weights are non-negative and sum to one, and $g_{\text{shape}}$ measures answer-shape violations such as missing required facets or insufficient distinct points. The deficit is a runtime state used to trigger verification or rescue; it is not a gold-label correctness score.
\end{definition}

\begin{definition}[Upgrade Utility]
The utility of escalation is the quality improvement obtained after spending additional retrieval or verification cost:
\begin{equation}
U_{\text{up}} = \Delta Q - \lambda \Delta \operatorname{cost},
\end{equation}
where $\Delta Q$ is measured only for offline analysis and is not available to the runtime policy. A central empirical question is whether a reference-free estimate of $D$ can predict positive $U_{\text{up}}$ rather than merely predict that a row is difficult.
\end{definition}

\textbf{Problem Statement.} Given $q$, $\mathcal{G}$, and a generated answer $a$, find: (1) a minimal proof subgraph $S^*$ that maximizes joint coverage; (2) the set of necessary evidence paths $N \subseteq S^*$; and (3) a submission decision $d^* \in \{\text{submit}, \text{monitor}, \text{block}, \text{abstain}\}$ that minimizes expected untrustworthiness.


% ══════════════════════════════════════════════════════════════
\section{Evidence Contract Construction}
\label{sec:mps}
% ══════════════════════════════════════════════════════════════

The evidence contract is instantiated by making the question obligations and answer claims explicit, then selecting a compact evidence set that covers both. MPS is therefore not an isolated graph algorithm in the system; it is the contract-construction primitive that produces the evidence state consumed by routing and submission.

\subsection{Algorithm}

The MPS optimizer solves a bi-objective weighted set cover:
\begin{equation}
\max_{S \subseteq \mathcal{E}} \alpha \cdot \operatorname{cov}_{\mathcal{O}}(S) + (1-\alpha) \cdot \operatorname{cov}_{\mathcal{C}}(S) \quad \text{s.t.} |S| \leq K, \operatorname{tokens}(S) \leq B.
\end{equation}

At each step, the greedy algorithm selects the evidence node with the highest marginal utility:
\begin{equation}
e^* = \arg\max_{e \in \mathcal{E} \setminus S} \frac{\Delta \operatorname{cov}_{\mathcal{O}}(e | S) + \lambda \cdot \Delta \operatorname{cov}_{\mathcal{C}}(e | S)}{\operatorname{tokens}(e)}.
\end{equation}

\begin{algorithm}[t]
\caption{Minimal Proof Subgraph Optimization}
\label{alg:mps}
\begin{algorithmic}[1]
\REQUIRE Evidence set $\mathcal{E}$, obligations $\mathcal{O}$, claims $\mathcal{C}$, budget $K$, $B$
\ENSURE Minimal proof subgraph $S^*$
\STATE $S \leftarrow \emptyset$
\WHILE{$|S| < K$ \AND $\operatorname{tokens}(S) < B$}
\STATE $e^* \leftarrow \arg\max_{e \in \mathcal{E} \setminus S} \frac{\Delta \operatorname{cov}_{\mathcal{O}}(e|S) + \lambda \cdot \Delta \operatorname{cov}_{\mathcal{C}}(e|S)}{\operatorname{tokens}(e)}$
\IF{$\Delta f(e^*|S) = 0$}
\STATE \textbf{break} \COMMENT{No marginal gain}
\ENDIF
\STATE $S \leftarrow S \cup \{e^*\}$
\ENDWHILE
\STATE \textit{// Minimality certificate: leave-one-out redundancy check}
\FOR{each $e_j \in S$}
\IF{$\operatorname{cov}_{\mathcal{O}}(S \setminus \{e_j\}) \geq \tau_O$ \AND $\operatorname{cov}_{\mathcal{C}}(S \setminus \{e_j\}) \geq \tau_C$}
\STATE $S \leftarrow S \setminus \{e_j\}$ \COMMENT{Prune redundant}
\ENDIF
\ENDFOR
\RETURN $S$
\end{algorithmic}
\end{algorithm}

\subsection{Theoretical Guarantee}

\begin{theorem}[Approximation Ratio]
\label{thm:mps}
Let $n = \max(P, M)$ where $P$ is the number of obligations and $M$ is the number of claims. The greedy algorithm (Algorithm~\ref{alg:mps}) achieves an approximation ratio of $O(\ln n)$ for the bi-objective weighted set cover problem.
\end{theorem}

\begin{proof}[Proof sketch]
The bi-objective function $f(S) = \alpha \cdot \operatorname{cov}_{\mathcal{O}}(S) + (1-\alpha) \cdot \operatorname{cov}_{\mathcal{C}}(S)$ is monotone submodular, as both coverage functions are submodular and a non-negative weighted sum of submodular functions is submodular. By the classical result of Nemhauser, Wolsey, and Fisher~\cite{nemhauser1978analysis}, the greedy algorithm for monotone submodular maximization under cardinality constraint achieves a $(1 - 1/e)$-approximation. For the weighted set cover formulation, the approximation ratio is $H_n = \sum_{i=1}^n 1/i = O(\ln n)$, where $H_n$ is the $n$-th harmonic number~\cite{chvatal1979greedy}.
\end{proof}


% ══════════════════════════════════════════════════════════════
\section{Deficit-Guided Escalation}
\label{sec:cpn}
% ══════════════════════════════════════════════════════════════

After the initial contract is built, the system must decide whether additional verification or graph rescue is worth its cost. CPN supplies the intervention primitive: it measures which evidence paths are necessary and whether a proposed escalation can replace a weak path without losing supported content. In the frozen SAGE implementation, the learned router estimates severe contract gaps, while the deterministic guard enforces preservation after repair.

\subsection{Intervention Framework}

For each $e_j \in S^*$, we perform a counterfactual intervention: remove $e_j$ and re-run the MPS optimizer on $\mathcal{E} \setminus \{e_j\}$. Let $S_{-j}$ denote the resulting subgraph. We measure:

\begin{itemize}
\item \textbf{Claim loss}: $\Delta_C(e_j) = \operatorname{cov}_{\mathcal{C}}(S^*) - \operatorname{cov}_{\mathcal{C}}(S_{-j})$
\item \textbf{Obligation loss}: $\Delta_O(e_j) = \operatorname{cov}_{\mathcal{O}}(S^*) - \operatorname{cov}_{\mathcal{O}}(S_{-j})$
\item \textbf{Sufficiency loss}: $\mathbb{1}[\text{proof\_sufficient}(S^*) \wedge \neg\text{proof\_sufficient}(S_{-j})]$
\item \textbf{Objective loss}: $\Delta_f(e_j) = f(S^*) - f(S_{-j})$
\end{itemize}

\subsection{Classification}

Based on the loss patterns, each evidence path is classified into one of five categories: \textit{claim-essential} ($\Delta_C > 0$), \textit{sufficiency-critical} (sufficiency loss), \textit{obligation-critical} ($\Delta_O > 0$), \textit{efficiency-critical} ($\Delta_f \geq 0.05$), or \textit{replaceable} (all losses negligible).

\begin{theorem}[Necessity Certification Correctness]
\label{thm:cpn}
Let $\epsilon$ be the coverage tolerance. The CPN certification has a false-positive rate (classifying replaceable evidence as necessary) of at most $\epsilon$, and a false-negative rate (classifying necessary evidence as replaceable) of $0$.
\end{theorem}

\begin{proof}[Proof sketch]
The false-negative rate is $0$ because the leave-one-out intervention exactly tests the necessity condition from Definition~3. The false-positive rate arises only when the re-optimization finds an alternative evidence path that provides equivalent coverage within tolerance $\epsilon$, which occurs with probability at most $\epsilon$ under the coverage gap assumption.
\end{proof}


% ══════════════════════════════════════════════════════════════
\section{Controlled Repair and Submission}
\label{sec:rcs}
% ══════════════════════════════════════════════════════════════

The final layer converts the contract state into an operational action: keep the text answer, escalate retrieval, accept a bounded repair, or abstain. RCS is the decision interface for this contract. The preservation guard is deliberately deterministic and reference-free at runtime; it can reject a repair that loses supported content even when the repair appears fluent.

\subsection{Risk Score Formulation}

Following Definition~\ref{def:risk}, the calibrated risk is:
\begin{align}
R(a, S) =\ & 0.34 \cdot r_{\text{audit}} + 0.26 \cdot r_{\text{unsupported}} + 0.18 \cdot r_{\text{proof\_gap}} \nonumber \\
& + 0.10 \cdot r_{\text{necessity}} + 0.07 \cdot r_{\text{contradiction}}.
\end{align}

The submission policy applies hierarchical thresholds:
\begin{itemize}
\item \textbf{Submit}: $R \leq 0.22$, $\operatorname{cov}_{\mathcal{C}} \geq 0.95$, $\operatorname{cov}_{\mathcal{O}} \geq 0.66$
\item \textbf{Submit with monitoring}: $R \leq 0.42$, $\operatorname{cov}_{\mathcal{C}} \geq 0.80$
\item \textbf{Block}: unsupported claims $> 0$ or contradicted claims $> 0$
\item \textbf{Defer/abstain}: otherwise
\end{itemize}

\subsection{PAC-Bayesian Generalization Bound}

\begin{theorem}[Generalization Bound]
\label{thm:rcs}
Let $\mathcal{D}$ be the distribution over question-evidence pairs, $m$ the sample size, and $\delta \in (0, 1)$. With probability at least $1 - \delta$ over the draw of $m$ i.i.d.\ samples, for any submission policy $\pi$ in the policy class $\Pi$:
\begin{equation}
\mathbb{E}_{(q, \mathcal{E}) \sim \mathcal{D}}[\text{Untrust}(\pi(q, \mathcal{E}))] \leq \hat{R}_{\text{emp}}(\pi) + \sqrt{\frac{\operatorname{KL}(\pi \| \pi_0) + \ln(2m/\delta)}{m}},
\end{equation}
where $\hat{R}_{\text{emp}}(\pi)$ is the empirical risk, $\operatorname{KL}(\pi \| \pi_0)$ is the KL divergence from the prior policy $\pi_0$, and $\text{Untrust}(\cdot)$ measures answer untrustworthiness.
\end{theorem}

\begin{proof}[Proof sketch]
This follows from the McAllester PAC-Bayesian theorem~\cite{mcallester1999pac}, applied to the policy class $\Pi$ of threshold-based submission rules parameterized by $(\tau_R, \tau_C, \tau_O)$. The calibrated risk score $R(a, S)$ provides a computable proxy for $\hat{R}_{\text{emp}}(\pi)$. The policy class complexity is bounded by $\operatorname{KL}(\pi \| \pi_0) \leq O(d \ln(m/d))$ where $d = 3$ is the number of threshold parameters.
\end{proof}


% ══════════════════════════════════════════════════════════════
\section{Experiments}
\label{sec:experiments}
% ══════════════════════════════════════════════════════════════

\subsection{Datasets and Protocol}

The current experimental package separates mechanism validation from frozen holdout evaluation. UltraDomain Mix130 is used for strict VSR ablation and candidate replay. Legal, Biology, and Domain-windowed slices provide external retrieval-baseline comparisons. Locked unseen75 and strict Repeat1--8 holdouts are source-line/source-context disjoint from prior development surfaces and are used only after freezing SAGE v33-preservation.

\begin{table}[htbp]
\caption{Current Evaluation Surfaces}
\label{tab:datasets}
\centering
\begin{tabular}{lrl}
\toprule
Surface & Questions & Role \\
\midrule
UltraDomain Mix130 & 130 & Strict VSR ablation and candidate replay \\
Legal/Biology/Domain external slices & 120 & Retrieval-baseline and semantic judge comparison \\
Locked unseen75 & 75 & Primary post-freeze zero-overlap holdout \\
Locked Repeat1--8 & 251 & Supplementary zero-overlap robustness and stress holdouts \\
\bottomrule
\end{tabular}
\end{table}

\subsection{Baselines}

We compare against \textbf{Naive RAG} and \textbf{HyDE} under same-context, same-metric settings. We also report \textbf{GraphRAG-lite} as a lightweight community graph baseline and \textbf{LightRAG current-130} as the evaluated representative heavy graph-RAG comparator. SAGE v33-preservation is treated as the frozen lead instantiation rather than a post-hoc tuned variant.

\subsection{Metrics}

We report lexical F1 and ROUGE-L, faithfulness, same-rubric LLM judge scores, pairwise A/B target win rates across five qualitative dimensions, runtime per question, and coverage/completion status. We keep lexical overlap and semantic judge results separate because they reward different answer properties.

\subsection{Results}

Tables~\ref{tab:strict-vsr-ablation-mix130}--\ref{tab:cost-latency-summary} are generated from completed artifacts by \texttt{experiments/build\_paper\_latex\_tables.py}; the corresponding Markdown/JSON package is stored under \texttt{results/paper\_tables/}.

% Auto-generated by experiments/build_paper_latex_tables.py
% Generated: 2026-09-18 20:34:32
% Source: results/paper_tables/paper_ready_experiment_tables.json
% Requires: \usepackage{booktabs}

\begin{table}[t]
\centering
\scriptsize
\caption{Strict VSR ablation on UltraDomain Mix130}
\label{tab:strict-vsr-ablation-mix130}
\begin{tabular}{lrrrrrrr}
\toprule
Config & Rows & F1 & ROUGE-L & Faith. & Repair & Abstain & s/q \\
\midrule
w/o\_all & 130/130 & 0.2276 & 0.1658 & 0.5000 & 0.0\% & 21.5\% & 9.0 \\
w/o\_evidence\_extraction & 130/130 & 0.2285 & 0.1612 & 0.5353 & 66.2\% & 40.0\% & 37.2 \\
w/o\_verification & 130/130 & 0.2138 & 0.1552 & 0.5000 & 0.0\% & 18.5\% & 17.6 \\
w/o\_repair & 130/130 & 0.2206 & 0.1591 & 0.8516 & 0.0\% & 23.1\% & 36.2 \\
full & 130/130 & 0.2437 & 0.1749 & 0.8516 & 34.6\% & 27.7\% & 44.2 \\
full\_v2 & 130/130 & 0.2824 & 0.2107 & 0.9728 & 40.0\% & 8.5\% & 41.0 \\
full\_v2\_rescue & 130/130 & 0.2878 & 0.2194 & 0.9756 & 27.7\% & 0.8\% & 41.7 \\
\bottomrule
\end{tabular}
\end{table}

\begin{table}[t]
\centering
\scriptsize
\caption{Paired tests for strict VSR ablations}
\label{tab:strict-vsr-ablation-tests}
\begin{tabular}{lrrrrr}
\toprule
Comparison & N & F1 Delta & F1 p & Faith. Delta & Faith. p \\
\midrule
full - w/o\_all & 130 & +0.0161 & 0.0277 & +0.3516 & <1e-15 \\
full - w/o\_evidence\_extraction & 130 & +0.0152 & 0.0641 & +0.3163 & 1.78e-15 \\
full - w/o\_verification & 130 & +0.0299 & 5.51e-07 & +0.3516 & <1e-15 \\
full - w/o\_repair & 130 & +0.0231 & 5.18e-05 & +0.0000 & 1.0000 \\
full - full\_v2 & 130 & -0.0386 & 2.51e-05 & -0.1211 & 2.03e-07 \\
full - full\_v2\_rescue & 130 & -0.0441 & 3.46e-05 & -0.1239 & 4.53e-07 \\
\bottomrule
\end{tabular}
\end{table}

\begin{table}[t]
\centering
\scriptsize
\caption{External retrieval baselines and SAGE v33}
\label{tab:retrieval-sage-v33}
\begin{tabular}{lrrrrrlr}
\toprule
Dataset & N & Naive & HyDE & VSR Rescue & SAGE v33 & Best & SAGE-Best \\
\midrule
Legal full context & 20 & 0.3632 & 0.3760 & 0.3496 & 0.3925 & hyde & +0.0164 \\
Biology full context & 50 & 0.2365 & 0.2385 & 0.2287 & 0.2385 & hyde & +0.0000 \\
Domain expansion windowed & 50 & 0.2711 & 0.2549 & 0.2337 & 0.2881 & naive & +0.0170 \\
\bottomrule
\end{tabular}
\end{table}

\begin{table}[t]
\centering
\scriptsize
\caption{Locked holdout weighted lexical results, Unseen75 plus Repeat1-8}
\label{tab:locked-holdout-weighted-lexical}
\begin{tabular}{lrrrrr}
\toprule
Method & N & F1 & ROUGE-L & s/q & Chars \\
\midrule
naive & 326 & 0.2972 & 0.2338 & 14.2 & 665.9 \\
hyde & 326 & 0.2972 & 0.2317 & 25.1 & 687.6 \\
sage\_v33\_preservation\_delta & 326 & 0.3164 & 0.2532 & 25.1 & 764.3 \\
\bottomrule
\end{tabular}
\end{table}

\begin{table}[t]
\centering
\scriptsize
\caption{Locked holdout weighted judge results, complete 326-row subset}
\label{tab:locked-holdout-weighted-judge}
\begin{tabular}{lrrrr}
\toprule
Method & N & Correct & Coverage & Overall \\
\midrule
naive & 326 & 3.092 & 2.856 & 2.960 \\
hyde & 326 & 3.110 & 2.856 & 2.942 \\
sage\_v33\_preservation\_delta & 326 & 3.816 & 3.469 & 3.555 \\
\bottomrule
\end{tabular}
\end{table}

\begin{table}[t]
\centering
\scriptsize
\caption{Locked holdout weighted pairwise win rates, complete 326-row subset}
\label{tab:locked-holdout-weighted-pairwise}
\begin{tabular}{lrr}
\toprule
Comparison & N & SAGE Target Win Rate \\
\midrule
sage\_v33\_vs\_hyde & 326 & 59.66\% \\
sage\_v33\_vs\_naive & 326 & 61.50\% \\
\bottomrule
\end{tabular}
\end{table}

\begin{table}[t]
\centering
\scriptsize
\caption{Repeat8 legal-only exhaustive stress holdout}
\label{tab:repeat8-legal54}
\begin{tabular}{lrrrrr}
\toprule
Method & N & F1 & ROUGE-L & Judge & s/q \\
\midrule
naive & 54 & 0.2958 & 0.2344 & 2.444 & 13.4 \\
hyde & 54 & 0.2862 & 0.2174 & 2.352 & 29.4 \\
sage\_v33\_preservation\_delta & 54 & 0.2862 & 0.2174 & 2.407 & 29.4 \\
\bottomrule
\end{tabular}
\vspace{0.35em}
\footnotesize{Repeat8 exhausts the remaining strict legal holdout capacity under the max-two-samples-per-context policy.}
\end{table}

\begin{table}[t]
\centering
\scriptsize
\caption{Cost and latency summary}
\label{tab:cost-latency-summary}
\begin{tabular}{lrrrrrrr}
\toprule
Dataset & N & Naive s/q & HyDE s/q & VSR s/q & VSR/Naive & VSR/HyDE & VSR-Best F1 \\
\midrule
Legal full context & 20 & 6.0 & 14.4 & 59.6 & 9.89 & 4.14 & -0.0264 \\
Biology full context & 50 & 6.8 & 20.7 & 46.9 & 6.88 & 2.26 & -0.0098 \\
Domain expansion windowed & 50 & 5.6 & 14.7 & 43.3 & 7.66 & 2.94 & -0.0374 \\
\bottomrule
\end{tabular}
\end{table}


\subsection{Interpretation}

The strict ablation supports the VSR mechanism claim: verification, repair, and rescue improve the internal no-repair baseline and materially reduce abstention on Mix130. The stronger paper-facing method is SAGE v33-preservation, which is reproducible, frozen before locked holdout evaluation, and safer than the exploratory v56 branch despite v56's aggregate replay gains.

Across locked unseen75 plus Repeat1--8, SAGE v33 improves weighted lexical F1 modestly and improves semantic judge quality substantially. Repeat8 is intentionally reported as a stress slice: it exhausts the remaining strict legal capacity, ties HyDE lexically, and trails Naive lexically, but it does not change the complete 326-row meta conclusion. We therefore frame the main claim as semantic-quality and preference improvement under strict holdouts, not as blanket lexical superiority over every retrieval baseline.

\subsection{Mechanism Depth Audit}

To test the central evidence-contract hypothesis directly, we aligned the frozen SAGE v33 Mix130 rows with the same-row `full\_v2\_rescue` outputs and computed a reference-free contract deficit from answer-shape and term-coverage signals. This is an observational audit rather than a randomized treatment comparison.

\begin{table}[htbp]
\caption{Evidence-Contract Mechanism Audit on Mix130}
\label{tab:contract-depth}
\centering
\begin{tabular}{lrrrr}
\toprule
Group & N & Mean F1 Delta & Wins/Losses & Mean Router Probability \\
\midrule
All aligned rows & 130 & +0.1768 & 113/17 & 0.2440 \\
Route skipped & 87 & +0.2089 & 83/4 & 0.0956 \\
Route selected & 43 & +0.1119 & 30/13 & 0.2991 \\
\bottomrule
\end{tabular}
\end{table}

The audit gives a useful negative result. Route-selected rows have larger estimated contract gaps and higher router probabilities, but they are also harder rows: their mean SAGE--VSR F1 delta is lower than for route-skipped rows. Router probability correlates negatively with the observed SAGE--VSR F1 delta ($r=-0.3782$), while the simple contract-gap proxy has a weaker negative correlation ($r=-0.1511$). Thus, the current contract mechanism is effective as a difficulty detector, but not yet as an upgrade-utility estimator. This separates two research problems that are often conflated: predicting that an answer is risky versus predicting that paying for escalation will improve it.

The next method iteration should therefore learn or calibrate an impact-aware escalation score, using only reference-free runtime features and validating it on a new locked holdout. The current results do not justify claiming that every high-deficit row benefits from escalation.


% ══════════════════════════════════════════════════════════════
\section{Related Work}
% ══════════════════════════════════════════════════════════════

\subsection{Graph-based RAG}
GraphRAG~\cite{edge2024graphrag} constructs community summaries from entity-relationship graphs for global retrieval. LightRAG~\cite{guo2024lightrag} employs dual-level retrieval combining entity and relationship matching. HippoRAG~\cite{gutierrez2024hipporag} uses personalized PageRank for multi-hop reasoning. None of these approaches formally verify whether the retrieved evidence constitutes a sufficient proof.

\subsection{Evidence Verification in RAG}
Self-RAG~\cite{asai2023self} introduces self-reflection tokens for retrieval gating. RAGAS~\cite{es2023ragas} provides automated evaluation metrics. ARES~\cite{yan2024ares} offers automated RAG evaluation. These approaches evaluate answer quality post-hoc but do not construct formal proof structures.

\subsection{Proof-based Verification}
EntailmentBank~\cite{dalvi2021entailment} generates multi-hop entailment trees. FOLIO~\cite{han2022folio} provides first-order logic reasoning. ProofWriter~\cite{tafjord2021proofwriter} performs natural language theorem proving. VSR-PathRAG is the first to combine formal proof structure with counterfactual necessity analysis and risk-calibrated submission gating.


% ══════════════════════════════════════════════════════════════
\section{Conclusion}
% ══════════════════════════════════════════════════════════════

We presented VSR-PathRAG as an evidence-contract-guided adaptive RAG framework. MPS, CPN, and RCS are unified as contract construction, deficit-guided escalation, and controlled submission primitives. The completed experiment package shows that the rescue-enhanced VSR layer improves the internal no-repair baseline, while the frozen SAGE v33-preservation instantiation improves semantic judge quality and pairwise preference on 326 strict locked-holdout rows. The mechanism audit also exposes the next scientific bottleneck: current contract-gap routing detects difficult rows, but its score is negatively associated with realized upgrade benefit. Future work should therefore target impact-aware escalation, official large-scale graph-baseline reruns, broader standard-benchmark validation, and human evaluation.


% ══════════════════════════════════════════════════════════════
\section*{Limitations}
% ══════════════════════════════════════════════════════════════

\begin{enumerate}
\item \textbf{LLM dependency}: Claim extraction and evidence signing rely on LLM quality.
\item \textbf{Scalability}: CPN's $O(|S| \times \text{MPS})$ complexity limits large $|S|$.
\item \textbf{Domain generalization}: The current strongest locked evidence comes from UltraDomain-derived domain slices; broader standard-benchmark validation remains future work.
\item \textbf{Baseline comparability}: GraphRAG-lite and LightRAG current-130 provide completed representative graph-RAG evidence, but full official graph-baseline reruns remain approximate or budget-dependent.
\item \textbf{Metric tension}: SAGE v33 is strong on same-rubric judge and pairwise preference, while lexical F1 gains over Naive RAG and HyDE are modest and not uniformly positive on every stress slice.
\item \textbf{Risk--benefit mismatch}: The current reference-free contract-gap router identifies difficult rows, but Mix130 analysis shows that high routing probability does not imply larger realized improvement after escalation.
\item \textbf{Theory}: The PAC-Bayesian bound assumes i.i.d.\ sampling, which may not hold for conversational RAG.
\end{enumerate}


% ══════════════════════════════════════════════════════════════
\begin{thebibliography}{99}
% ══════════════════════════════════════════════════════════════

\bibitem{edge2024graphrag}
D.~Edge et al., ``From local to global: A graph rag approach to query-focused summarization,'' \textit{arXiv preprint arXiv:2404.16130}, 2024.

\bibitem{guo2024lightrag}
Z.~Guo et al., ``Lightrag: Simple and fast graph-based retrieval-augmented generation,'' \textit{arXiv preprint arXiv:2410.05779}, 2024.

\bibitem{gutierrez2024hipporag}
B.~J. Guti\'errez et al., ``Hipporag: Neurobiologically inspired long-term memory for large language models,'' \textit{NeurIPS}, 2024.

\bibitem{asai2023self}
A.~Asai et al., ``Self-rag: Learning to retrieve, generate, and critique through self-reflection,'' \textit{ICLR}, 2024.

\bibitem{es2023ragas}
S.~Es et al., ``Ragas: Automated evaluation of retrieval augmented generation,'' \textit{arXiv preprint arXiv:2309.15217}, 2023.

\bibitem{yan2024ares}
S.~Yan et al., ``Ares: Automated rag evaluation system,'' \textit{NAACL}, 2024.

\bibitem{dalvi2021entailment}
B.~Dalvi et al., ``Explaining answers with entailment trees,'' \textit{EMNLP}, 2021.

\bibitem{han2022folio}
S.~Han et al., ``Follio: Natural language reasoning with first-order logic,'' \textit{arXiv preprint arXiv:2209.00840}, 2022.

\bibitem{tafjord2021proofwriter}
O.~Tafjord et al., ``Proofwriter: Generating implications, proofs, and abductive statements over natural language,'' \textit{ACL Workshop}, 2021.

\bibitem{nemhauser1978analysis}
G.~L. Nemhauser, L.~A. Wolsey, and M.~L. Fisher, ``An analysis of approximations for maximizing submodular set functions--I,'' \textit{Mathematical Programming}, vol.~14, pp.~265--294, 1978.

\bibitem{chvatal1979greedy}
V.~Chv\'atal, ``A greedy heuristic for the set-covering problem,'' \textit{Mathematics of Operations Research}, vol.~4, no.~3, pp.~233--235, 1979.

\bibitem{mcallester1999pac}
D.~McAllester, ``Pac-bayesian model averaging,'' in \textit{COLT}, 1999, pp.~164--170.

\bibitem{lewis2020retrieval}
P.~Lewis et al., ``Retrieval-augmented generation for knowledge-intensive nlp tasks,'' \textit{NeurIPS}, 2020.

\bibitem{gao2022hallucination}
L.~Gao et al., ``The curse of hallucination in neural text generation,'' \textit{ACM Computing Surveys}, 2022.

\bibitem{wei2022chain}
J.~Wei et al., ``Chain-of-thought prompting elicits reasoning in large language models,'' \textit{NeurIPS}, 2022.

\end{thebibliography}

\end{document}
